\chapter{Cálculo de Malliavin y rough paths}\label{cap:rough-paths}

\begin{panorama}
La teoría de rough paths de Lyons (1998) resuelve ecuaciones diferenciales dirigidas por
señales irregulares \emph{trayectoria a trayectoria}, sin martingalas. Para procesos gaussianos, la
solución es entonces un funcional continuo de un \emph{levantamiento} del ruido, y el cálculo
de Malliavin se puede hacer en sentido fuerte: la solución es Fréchet-diferenciable en las
direcciones de Cameron--Martin. Presentamos los ingredientes básicos, el teorema de
Hörmander para EDRs gaussianas (Cass--Friz, Cass--Hairer--Litterer--Tindel) y algunos
desarrollos recientes.
\end{panorama}

\section{Rough paths en pocas palabras}
Sea $\alpha\in(\frac13,\frac12]$. Un \emph{rough path} $\alpha$-Hölder en $\R^m$ es un par $\mathbf X=(X,\mathbb X)$ con $X:[0,T]\to\R^m$ y
$\mathbb X:[0,T]^2\to\R^{m\times m}$ tales que
\[
  \abs{X_t-X_s}\le C\abs{t-s}^\alpha,\qquad\abs{\mathbb X_{s,t}}\le C\abs{t-s}^{2\alpha},\qquad\mathbb X_{s,t}-\mathbb X_{s,u}-\mathbb X_{u,t}=X_{s,u}\otimes X_{u,t}
\]
(\emph{relación de Chen}), con $X_{s,t}=X_t-X_s$. Heurísticamente, $\mathbb X_{s,t}=\int_s^tX_{s,r}\otimes\dd X_r$: es la información
de segundo orden que la trayectoria por sí sola no determina cuando $\alpha\le\frac12$.

\begin{lema}[Lema de cosido, \textup{sewing lemma}]\label{lem:sewing}
Sea $\Xi_{s,t}$ una función de dos parámetros con
$\abs{\Xi_{s,t}-\Xi_{s,u}-\Xi_{u,t}}\le C\abs{t-s}^\theta$ para $s\le u\le t$, con $\theta>1$. Entonces existe una única función
aditiva $\mathcal I_t$ con $\abs{\mathcal I_t-\mathcal I_s-\Xi_{s,t}}\le C'\abs{t-s}^\theta$, dada por $\mathcal I_t-\mathcal I_s=\lim_{\abs\pi\to0}\sum_{[u,v]\in\pi}\Xi_{u,v}$.
\end{lema}

\begin{proof}[Esbozo]
Para particiones diádicas, quitar un punto de una partición de $n$ intervalos cambia
la suma en a lo sumo $C(\frac{2\abs{t-s}}{n})^\theta$. Eliminando puntos sucesivamente hasta llegar a $\{s,t\}$, la
diferencia total está acotada por $C\abs{t-s}^\theta\sum_nn^{-\theta}<\infty$, porque $\theta>1$. Esto da existencia del
límite y la estimación; la unicidad se sigue de que una función aditiva con $\abs{f_{s,t}}\le C\abs{t-s}^\theta$,
$\theta>1$, es nula.
\end{proof}

Con $\Xi_{s,t}=Y_sX_{s,t}+Y_s'\mathbb X_{s,t}$, para $Y$ \emph{controlado} por $X$ (es decir, $Y_{s,t}=Y_s'X_{s,t}+O(\abs{t-s}^{2\alpha})$),
se tiene $\theta=3\alpha>1$ y el lema define la \emph{integral rugosa} $\int Y\,\dd\mathbf X$. Esto permite resolver
\begin{equation}\label{eq:edr}
  \dd Y_t=V_0(Y_t)\,\dd t+V(Y_t)\,\dd\mathbf X_t
\end{equation}
por punto fijo, y la solución depende de manera \emph{localmente Lipschitz} de $\mathbf X$ en la métrica de
rough paths (\emph{teorema de continuidad de Lyons}, o teorema del límite universal).

\begin{ejemplo}[Browniano]
$\mathbb X^{\text{Itô}}_{s,t}=\int_s^tW_{s,r}\otimes\dd W_r$ y $\mathbb X^{\text{Strat}}_{s,t}=\mathbb X^{\text{Itô}}_{s,t}+\frac12(t-s)I$ son dos rough
paths distintos sobre la misma trayectoria; con el primero,~\eqref{eq:edr} es la EDE de Itô y, con el
segundo, la de Stratonovich. Ambos existen c.s.\ (por Kolmogorov en dos parámetros) para todo
$\alpha<\frac12$.
\end{ejemplo}

\section{Levantamientos gaussianos}

\begin{teorema}[Friz--Victoir]
Sea $X=(X^1,\dots,X^m)$ un proceso gaussiano centrado
continuo con componentes independientes, cuya covarianza tiene $\rho$-variación bidimensional
finita con $\rho<2$. Entonces $X$ se levanta canónicamente a un rough path geométrico $\mathbf X$ (límite
de los levantamientos de aproximaciones lineales a trozos) de $p$-variación finita para todo
$p>2\rho$. Para el browniano fraccionario, $\rho=\frac1{2H}$, de modo que la condición es $H>\frac14$.
\end{teorema}

Para $H\le1/4$, las aproximaciones lineales a trozos de $\int B^{H,1}\,\dd B^{H,2}$ divergen (la varianza del
área de Lévy explota), en consonancia con la fórmula de Itô para la divergencia del Capítulo~\ref{cap:fraccionario}.

\section{Diferenciabilidad de Malliavin trayectorial}

\begin{teorema}
Sea $X$ gaussiano como antes, con espacio de Cameron--Martin $\Hc$, y
supongamos que $\Hc$ se inyecta en el espacio de trayectorias de $q$-variación finita con $\frac1p+\frac1q>1$
\textup{(}regularidad complementaria de Young\textup{)}. Sea $Y=\Psi(\mathbf X)$ la solución de~\eqref{eq:edr} con $V\in C_b^\infty$.
Entonces, para cada trayectoria, $h\mapsto\Psi(T_h\mathbf X)$ es Fréchet-diferenciable en $\Hc$, donde $T_h\mathbf X$ es
el rough path trasladado, y
\[
  \ip{DY_t}{h}_\Hc=\int_0^tJ_{t\leftarrow s}V(Y_s)\,\dd h_s,
\]
con $J_{t\leftarrow s}$ el jacobiano del flujo, solución de la ecuación linealizada. En particular, $Y_t\in\D^{1,2}_{\mathrm{loc}}$
y, si $X$ es el fBm con $H>1/4$, la condición de Young complementaria se satisface.
\end{teorema}

Es la misma fórmula que en el Teorema~\ref{teo:derivada-ede}, pero ahora con una integral de Young respecto
a $h$ y válida $\omega$ a $\omega$. Para el fBm, $\Hc\subset C^{q\text{-var}}$ para todo $q>(H+\frac12)^{-1}$, y $\frac1p+\frac1q>1$ se reduce a
$H>\frac14$.

\section{El teorema de Hörmander para EDRs gaussianas}

\begin{teorema}[Cass--Friz 2010; Cass--Hairer--Litterer--Tindel 2015]
Sea $X$ un proceso
gaussiano que satisface las hipótesis anteriores más una condición de no degeneración \textup{(}no
determinismo local\textup{)}, por ejemplo el fBm con $H\in(\frac14,\frac12]$. Si los campos $V_0,V_1,\dots,V_m$
satisfacen la condición de Hörmander en $y_0$, entonces para todo $t>0$:
\begin{enumerate}[label=(\roman*)]
  \item la ley de $Y_t$ es absolutamente continua (Cass--Friz~\cite{CassFriz10});
  \item su densidad es $C^\infty$ (Cass--Hairer--Litterer--Tindel~\cite{CassHairerLittererTindel15}).
\end{enumerate}
\end{teorema}

\begin{proof}[Ideas de la prueba]
Hay tres obstáculos respecto al caso browniano, y cada uno requirió una
idea nueva.

\emph{(a) Sustituto del lema de Norris.} No hay martingalas, así que el argumento probabilístico
del Lema~\ref{lem:norris} no sirve. Se sustituye por un \emph{lema de Norris determinista} basado en la \emph{rugosidad
verdadera} (Hairer--Pillai): una trayectoria $X$ es $\theta$-rugosa si, en todo intervalo y en toda dirección,
no puede ser ``demasiado regular''. Para trayectorias rugosas, si $\int\phi\,\dd X+\int\psi\,\dd t$ es pequeño,
entonces $\phi$ y $\psi$ son pequeños (unicidad cuantitativa de la descomposición tipo Doob--Meyer).
La rugosidad se verifica c.s.\ para procesos gaussianos con no determinismo local.

\emph{(b) Integrabilidad del jacobiano.} Para una EDE, $J$ tiene todos los momentos por BDG--Gronwall.
Para una EDR, la cota natural es $\abs J\le\exp(C\norm{\mathbf X}^p_{p\text{-var}})$, y $\norm{\mathbf X}^p$ con $p>2$ \emph{no}
tiene colas gaussianas: la exponencial no es integrable. Cass, Litterer y Lyons resolvieron este
problema con las \emph{particiones voraces}: el jacobiano está controlado por $\exp(CN(\mathbf X))$, donde $N$
es el número de intervalos en que la $p$-variación del rough path supera un umbral, y $N$ sí tiene
colas gaussianas por el teorema de Borell aplicado a la traslación de Cameron--Martin.

\emph{(c) La matriz de Malliavin.} La norma de $\Hc$ no es local en tiempo, de modo que $\ip{DY}{DY}_\Hc$
no se escribe como una integral en $t$. Se necesitan \emph{desigualdades de interpolación} que acoten
inferiormente $\norm f_\Hc$ en términos de normas locales de $f$; para el fBm con $H<1/2$, $\norm f_\Hc\gtrsim\norm f_{L^2}$
en cierto sentido, lo que permite reducir al argumento de los corchetes.
\end{proof}

\section{Desarrollos recientes}
\begin{itemize}
  \item \textbf{Ecuaciones diferenciales estocásticas rugosas.} Friz, Hocquet y Lê (2021) desarrollaron
  una teoría para ecuaciones dirigidas simultáneamente por un browniano (tratado
  con Itô) y un rough path (tratado trayectorialmente), $\dd Y=b\,\dd t+\sigma\,\dd W+f\,\dd\mathbf X$. Son
  el lenguaje natural para filtrado no lineal robusto, volatilidad local-estocástica y control
  estocástico con ruido común. Su análisis de Malliavin combina ambos mundos.
  \item \textbf{TCL trayectoriales y desviaciones moderadas.} Se han obtenido TCLs y principios
  de desviaciones moderadas para EDPEs con ruido multiplicativo vía rough paths, donde
  el cálculo de Malliavin pathwise permite linealizar alrededor del límite determinista.
  \item \textbf{Densidades en régimen rugoso.} Inahama~\cite{Inahama14} y otros estudiaron la diferenciabilidad
  de Malliavin y los desarrollos asintóticos (Laplace, Watanabe) de soluciones de EDRs, con
  aplicaciones a núcleos de calor en geometría sub-riemanniana.
  \item \textbf{Signaturas.} La signatura de un rough path, la sucesión de todas sus integrales iteradas,
  determina la trayectoria salvo reparametrizaciones (Hambly--Lyons; Boedihardjo, Geng,
  Lyons, Yang). Es la base de métodos de aprendizaje automático para series temporales, y
  su esperanza bajo la medida de Wiener se calcula con las mismas herramientas de caos de
  Wiener del Capítulo~\ref{cap:caos}.
\end{itemize}
